% 3-3-3-eda.tex
%  Budget: 3pp.
%  Source map: seasonality/volatility/spike characteristics of electricity
%  price -> jiang_review_2018, ugurlu_electricity_2018, alkawaz_day-ahead_2022.
%  Figures 01-09 already exported to reports/figures/ (HANDOFF §19.2).
%  Numbers: price stats from 3-3-1 (same source). 84.04 appears ONCE, in its
%  permitted DESCRIPTIVE sense only -- it is NOT the regime threshold; the
%  threshold is 62.6989 and is defined in 3-8.
%  Rules: \lr{} every numeral; \lr{—} for em-dash.

\subsection{تحلیل اکتشافی داده}

پیش از مدل‌سازی، ویژگی‌های آماری سری قیمت مستند شده‌اند. هدف این بخش
توصیف است و نه استنتاج؛ هر ادعای علّی درباره سازوکار قیمت‌گذاری خارج از
دامنه این پژوهش است.

\subsubsection*{توزیع قیمت}

توزیع قیمت (شکل \lr{1}) نامتقارن است و دنباله راست بلندتری دارد. مقادیر
منفی، چنان که در بخش \lr{3-3-1} آمد، \lr{538} ساعت را در بر می‌گیرند.
وجود این مقادیر، سه پیامد روش‌شناختی دارد: کنارگذاشتن معیارهای نسبی
مبتنی بر تقسیم (بخش \lr{3-5})، ناکارآمدی تبدیل لگاریتمی که در بسیاری از
مطالعات سری‌های مالی متعارف است، و لزوم پرهیز از هر پیش‌فرض مثبت بودن در
پیاده‌سازی مدل‌ها.

\subsubsection*{فصلی‌بودن چندگانه}

سه دوره تناوب هم‌زمان در داده دیده می‌شود و هر سه در ادبیات این حوزه
شناخته‌شده‌اند \cite{jiang_review_2018}:

\begin{itemize}
  \item \textbf{روزانه} (شکل \lr{2}) \lr{—} الگوی درون‌روزی با دو قله
        صبح و عصر، که بازتاب منحنی بار است.
  \item \textbf{هفتگی} (شکل \lr{3}) \lr{—} تفاوت روشن میان روزهای کاری و
        آخر هفته. این الگو، دلیل حضور متغیرهای مجازی روز هفته در مجموعه
        ویژگی‌ها (بخش \lr{3-4}) است.
  \item \textbf{سالانه} (شکل \lr{4}) \lr{—} تغییرات فصلی ناشی از تقاضای
        گرمایش و سرمایش و از الگوی تولید تجدیدپذیر.
\end{itemize}

حضور هم‌زمان این سه دوره تناوب، یکی از دلایلی است که مدل‌های سری زمانی
تک‌فصلی برای این مسئله ناکافی‌اند.

\subsubsection*{خوشه‌بندی نوسان و مقادیر حدی}

شکل \lr{5} خوشه‌بندی نوسان را نشان می‌دهد: دوره‌های پرنوسان و کم‌نوسان به
جای پراکندگی یکنواخت، در کنار هم تجمع می‌کنند. این ویژگی، همراه با
جهش‌های ناگهانی و کوتاه‌مدت، از خصوصیات شناخته‌شده قیمت برق است
\cite{ugurlu_electricity_2018, alkawaz_day-ahead_2022}.

اهمیت این مشاهده برای طرح پژوهش آن است که خطای مدل در سراسر داده یکنواخت
توزیع نمی‌شود. مدلی که بر میانگین خطا بهینه شده، در ساعات پرتنش \lr{—}
یعنی جایی که خطا پرهزینه‌ترین است \lr{—} ضعیف‌ترین عملکرد را دارد. این
ناهمگنی، انگیزه مستقیم وزن‌دهی مشروط به رژیم در بخش \lr{3-8} است.

در همین چارچوب، یک آماره توصیفی درخور ذکر است: آستانه سه‌سیگما بر داده
آموزش برابر \lr{84.04} یورو بر مگاوات‌ساعت است. این عدد در این پژوهش
\emph{تنها} یک آماره توصیفی است و آستانه تفکیک رژیم نیست؛ آستانه عملیاتی
مورد استفاده \lr{62.6989} است و دلیل انتخاب آن در بخش \lr{3-8} می‌آید.

\subsubsection*{ساختار خودهمبستگی}

نمودارهای خودهمبستگی و خودهمبستگی جزئی (شکل \lr{6}) قله‌های روشنی در
وقفه‌های \lr{24} و \lr{168} ساعت نشان می‌دهند؛ یعنی همان دوره‌های روزانه
و هفتگی. این ساختار مستقیماً تعیین می‌کند کدام وقفه‌ها وارد مجموعه
ویژگی‌ها شوند: روزهای \lr{1}، \lr{2}، \lr{3} و \lr{7} پیش از روز هدف
(بخش \lr{3-4}).

\subsubsection*{متغیرهای برون‌زا}

شکل \lr{7} همبستگی میان قیمت و دو متغیر برون‌زا را نشان می‌دهد: همبستگی
مثبت با پیش‌بینی بار و همبستگی منفی با پیش‌بینی تولید تجدیدپذیر. جهت هر
دو با انتظار فیزیکی سازگار است؛ تولید با هزینه نهایی نزدیک به صفر،
منحنی عرضه را جابه‌جا می‌کند و قیمت تسویه را پایین می‌آورد.

شکل \lr{7b} همین همبستگی‌ها را به‌صورت غلتان در طول زمان نشان می‌دهد و
نکته مهم‌تری را آشکار می‌کند: این همبستگی‌ها ثابت نیستند. قدرت رابطه میان
قیمت و تولید تجدیدپذیر در طول دوره تغییر می‌کند، که با گسترش ظرفیت نصب‌شده
در همان سال‌ها سازگار است.

\subsubsection*{شکست‌های ساختاری و سری روزانه}

شکل \lr{8} کل سری را همراه با نقاط شکست ساختاری نشان می‌دهد، و شکل
\lr{9} سری میانگین بار پایه روزانه را که هدف دوم این پژوهش است.

مشاهده مربوط به شکست‌های ساختاری، به فرض یکم (ایستایی) در بخش \lr{3-2}
بازمی‌گردد و آن را تعدیل می‌کند: سری قیمت در سراسر دوره ایستا نیست. پاسخ
عملی این پژوهش به آن، واسنجی مجدد روزانه بر پنجره‌ای غلتان است و نه
آموزش یک‌باره بر کل تاریخچه.
